% BHU PYQ -- Professional Exam Template
\documentclass[11pt,a4paper]{article}

\usepackage[a4paper,top=2.5cm,bottom=2.5cm,left=2.8cm,right=2.8cm,headheight=14pt]{geometry}
\usepackage{amsmath,amssymb}
\usepackage[T1]{fontenc}
\usepackage[utf8]{inputenc}
\usepackage{lmodern}
\usepackage{microtype}
\usepackage{fancyhdr}
\usepackage{enumitem}
\usepackage{hyperref}
\usepackage{lastpage}
\usepackage{parskip}
\usepackage{tikz}
\usetikzlibrary{arrows.meta,shapes,positioning}

\hypersetup{colorlinks=true,urlcolor=black,linkcolor=black,
  pdftitle={B.Sc. (Hons.) Zoology Sem-VI ZOB-602 -- BHU PYQ}}

\pagestyle{fancy}
\fancyhf{}
\renewcommand{\headrulewidth}{0.4pt}
\renewcommand{\footrulewidth}{0.4pt}
\fancyhead[L]{\small   \textit{Banaras Hindu University}}
\fancyhead[R]{\small   \textit{Zoology Paper: ZOB-602}}
\fancyfoot[C]{\small Page  \thepage\ of \pageref{LastPage}}


\newlist{parts}{enumerate}{3}
\setlist[parts,1]{label=(\alph*),leftmargin=*,itemsep=4pt,topsep=3pt}
\setlist[parts,2]{label=(\roman*),leftmargin=*,itemsep=2pt,topsep=2pt}
\setlist[parts,3]{label=(\arabic*),leftmargin=*,itemsep=2pt,topsep=2pt}

\newcommand{\pts}[1]{\hfill   \textit{[#1]}}


\begin{document}


\begin{center}
  {\large\bfseries BANARAS HINDU UNIVERSITY}\\[2pt]
  {
\normalsize B.Sc. (Hons.) Semester VI Examination, 2013-14}\\[6pt]
  \rule{0.8\linewidth}{0.5pt}\\[6pt]
  {\large\bfseries Zoology}\\[4pt]
  {\large\bfseries Paper: ZOB-602 --- Cell Biology, Genetics and Evolution}\\[4pt]
  {
\normalsize Time: Three hours \hfill Full Marks: 70}
\end{center}

\rule{\linewidth}{0.4pt}

\smallskip



\noindent   \textit{Note: Answer two questions from each Section including Question No. 1 of each Section which is compulsory. Use a separate answer book for each Section. The figures in the right-hand margin indicate marks.}

\bigskip

\begin{center}
   \textbf{SECTION -- A} \\
   \textbf{(Cell Biology) (Marks: 26)}
\end{center}

\smallskip



\noindent   \textit{Note: Answer two questions including Question No. 1 which is compulsory.}

\medskip


\noindent   \textbf{Question 1.} Answer the following parts:

\begin{parts}
  \item Write whether the following statements are True or False, giving reasons for your answer in 3--4 lines:\pts{$2  \times 2 = 4$}
  
\begin{parts}
    \item The endoplasmic reticulum lumen is full of soluble Luminal resident proteins and enzymes, whereas the resident proteins in the Golgi apparatus are all membrane bound.
    \item Lysosome contains only hydrolases.
  \end{parts}

  \item Select the most suitable answer for the following:\pts{$2  \times \frac{1}{2} = 1$}
  
\begin{parts}
    \item Multivesicular Bodies (MVBs) are found in:
    
\begin{parts}
      \item early endosomes
      \item late endosomes
      \item hybrid endosomes
      \item lysosomes
    \end{parts}

    \item KDEL protein is retained in:
    
\begin{parts}
      \item cytoplasm
      \item nucleus
      \item endoplasmic reticulum
      \item Golgi apparatus
    \end{parts}
  \end{parts}

  \item Differentiate between the following:\pts{$2  \times 1\frac{1}{2} = 3$}
  
\begin{parts}
    \item Retrograde transport and Anterograde transport
    \item Active transport and Passive transport.
  \end{parts}
\end{parts}

\medskip


\noindent   \textbf{Question 2.} Answer the following parts:

\begin{parts}
  \item Define signal peptides. Explain entry and passage of proteins through ER.\pts{$2 + 7 = 9$}
  \item Explain the events in different phases of cell cycle. Describe the genetic regulation of the cell cycle.\pts{$3 + 6 = 9$}
\end{parts}

\medskip


\noindent   \textbf{Question 3.} Answer the following parts:

\begin{parts}
  \item Describe the structure of Lampbrush chromosome and explain its significance in the study of gene expression.\pts{$4 + 5 = 9$}
  \item Write short notes on the following:\pts{$3  \times 3 = 9$}
  
\begin{parts}
    \item Centromere
    \item Nuclear pore complex
    \item Endosomes.
  \end{parts}
\end{parts}

\pagebreak

\begin{center}
   \textbf{SECTION -- B} \\
   \textbf{(Genetics) (Marks: 26)}
\end{center}

\smallskip



\noindent   \textit{Note: Answer two questions including Question No. 1 which is compulsory.}

\medskip


\noindent   \textbf{Question 1.} Answer the following parts:

\begin{parts}
  \item Write whether the following statements are True or False giving valid reasons for your answer:\pts{$2  \times 2 = 4$}
  
\begin{parts}
    \item The on/off switch of   \textit{sxl} gene of Drosophila is not sensitive to the X : A ratio.
    \item Frameshift mutations are more detrimental than base substitution mutations.
  \end{parts}

  \item The promoter region of an operon has a mutation to which RNA polymerase cannot bind and initiate transcription. Will this mutation prevent the binding of RNA polymerase to another copy of operon in trans condition? Explain with reasons.\pts{2}

  \item Select the best options from the choices given below:\pts{$4  \times \frac{1}{2} = 2$}
  
\begin{parts}
    \item Out of 800 progeny of a three-point testcross there were 12 double crossover recombinants, whereas 20 had been expected on the basis of no interference. The interference is:
    
\begin{parts}
      \item 0.10
      \item 0.20
      \item 0.30
      \item 0.40
    \end{parts}

    \item What will be the state of lac operon of a   \textit{lac I} mutant?
    
\begin{parts}
      \item The   \textit{lac} genes would be expressed efficiently only in the absence of lactose.
      \item The   \textit{lac} genes would be expressed efficiently only in the presence of lactose.
      \item The   \textit{lac} genes would be expressed continuously.
      \item The   \textit{lac} genes would never be expressed efficiently.
    \end{parts}

    \item Phenylketonuria (PKU) arises because of:
    
\begin{parts}
      \item inclusion bodies in the neurons
      \item metabolic dsyfunction
      \item multifactorial causes
      \item transposable element
    \end{parts}

    \item A strain of   \textit{E. coli} is Lac$^-$ and Met$^-$. After transformation experiment to select only for Met$^+$ cells, it should be plated on:
    
\begin{parts}
      \item minimal media + glucose + methionine
      \item minimal media + glucose
      \item rich media
      \item minimal media + lactose + methionine
    \end{parts}
  \end{parts}
\end{parts}

\medskip


\noindent   \textbf{Question 2.} Answer the following parts:

\begin{parts}
  \item Amino acid glycine is encoded by CGA. A single base substitution changes it to a nonsense codon. Specify the base change that would have taken place to convert it into a nonsense codon. Is this change a transversion or transition? Distinguish between the two.\pts{$1 + 1 + 2 = 4$}
  \item Two-point testcrosses revealed the following map results:
  
\begin{parts}
    \item A --- B 17 map units
    \item B --- C 11 map units.
  \end{parts}
  What are the two possible maps involving the three given loci? Which additional cross should be performed to resolve the correct map?\pts{$1 + 3 = 4$}
  \item An F' (F prime) element containing the maltose gene (   \textit{mal}) is found in   \textit{E. coli}. Explain in stepwise manner how this element would have been produced starting from an F$^+$ plasmid.\pts{6}
  \item What are DNA arrays? How are they useful in studying gene expression?\pts{4}
\end{parts}

\medskip


\noindent   \textbf{Question 3.} Answer the following parts:

\begin{parts}
  \item The height of a plant is controlled by two pairs of alleles assorting independently, A, a and B, b. Each dominant allele, A and B contribute 40 cm to the height and each recessive alleles a and b contribute 10 cm to the height. Assuming that these loci have additive effects, what will be the height of the F$_1$ progeny from a cross of AABB $ \times$ aabb? How many phenotypic classes are possible in F$_2$ progeny? And also mention the range of height of F$_2$ progeny.\pts{$2 + 3 = 5$}
  \item Describe the three types of tetrad that can be obtained from tetrad analysis of Neurospora. How can one distinguish whether the two genes are linked or on different chromosomes on the basis of type of the tetrads observed?\pts{$2 + 3 = 5$}
  \item A transposable element has a deletion in the transposase gene. What will be the consequence of this deletion? Under what conditions can this transposable element jump in the genome?\pts{3}
  \item Briefly write on triplet repeat expansion disorder taking Huntington's disease as an example. How is it different from a multifactorial disease?\pts{$3 + 2 = 5$}
\end{parts}

\pagebreak

\begin{center}
   \textbf{SECTION -- C} \\
   \textbf{(Evolution) (Marks: 18)}
\end{center}

\smallskip



\noindent   \textit{Note: Answer two questions including Question No. 1 which is compulsory.}

\medskip


\noindent   \textbf{Question 1.} Answer the following parts:

\begin{parts}
  \item State whether the following statements are True or False, giving reasons in 2--3 sentences:\pts{$2  \times 2 = 4$}
  
\begin{parts}
    \item Toothed birds originated in Cambrian period and became extinct by the end of Cretaceous period.
    \item Eohippus possessed 3 toes in fore- and hind-limbs.
  \end{parts}

  \item Define the following:\pts{$2  \times 1 = 2$}
  
\begin{parts}
    \item Mullerian mimicry
    \item Hardy Weinberg law of equilibrium.
  \end{parts}
\end{parts}

\medskip


\noindent   \textbf{Question 2.} Answer the following parts:

\begin{parts}
  \item What is molecular phylogeny? How can DNA sequences or proteins be used to establish phylogenetic relationships among the different group of animals?\pts{6}
  \item Explain the following:\pts{$2  \times 3 = 6$}
  
\begin{parts}
    \item Radioactive dating
    \item Adoptive colouration.
  \end{parts}
\end{parts}

\medskip


\noindent   \textbf{Question 3.} Answer the following parts:

\begin{parts}
  \item Describe the various means of prezygotic reproductive isolation.\pts{6}
  \item Write short notes on the following:\pts{$2  \times 3 = 6$}
  
\begin{parts}
    \item Allopatric speciation
    \item Chromosomal polymorphism.
  \end{parts}
\end{parts}

\vfill
\rule{\linewidth}{0.4pt}

\begin{center}\small
\href{https://bhu-exam.vercel.app/}{Click Here}\\[4pt]\href{https://me-aryan.vercel.app/}{Click Here}\\[4pt]\href{https://quashan-me.vercel.app/}{Click Here}
\end{center}

\end{document}
